\clearpage\begingroup\NoSourceHead\SourceLeaf{697}
\begin{center}\bfseries
BEMERKUNGEN\\[3pt]
ZUR ZWEITEN ABTEILUNG DES ZEHNTEN BANDES.
\end{center}
\noindent\hspace*{\parindent}Im Jahre 1911 entwickelte \textsc{F. Klein} den Plan, der Gesamtausgabe der \textsc{Gauss}schen Werke eine
ausführliche wissenschaftliche Biographie von \textsc{Gauss} anzufügen. (Vergl. die Angaben in den Bemerkungen
zu Werke XI, 2.) Das dort über die in jenem Bande vereinigten drei Abhandlungen — »Essays« — Gesagte
gilt mutatis mutandis auch von den sieben Abhandlungen des vorliegenden Bandes. Mit diesen zehn Essays
erscheint das \textsc{Klein}sche Programm der wissenschaftlichen Biographie von \textsc{Gauss} durchgeführt. Da diese
10 Essays das wissenschaftliche Lebenswerk von \textsc{Gauss}, im besonderen also auch die von \textsc{Gauss} selbst
nicht veröffentlichten Untersuchungen, soweit sie im Nachlass und im Briefwechsel niedergelegt sind, zusammenfassend
darstellen und erläutern, dienen sie über ihre unmittelbare biographisch-historische Zielsetzung
hinaus auch der Bearbeitung des Nachlasses, den sie vielfach ergänzen und vervollständigen. In
bezug auf das letztere mag es genügen, darauf hinzuweisen, dass \textsc{Gauss}’ populärste Erfindung, der elektromagnetische
Telegraph, in den 12 Bänden, die die eigentlichen gesammelten Werke von \textsc{Gauss} enthalten,
nicht ausführlich beschrieben ist; nur der auf diese Erfindung bezügliche Briefwechsel wird im Band XI, 1
wiedergegeben, während die eingehende Beschreibung des von \textsc{Gauss} angegebenen und konstruierten
Apparates, als nicht von \textsc{Gauss} selbst herrührend, sich in dem physikalischen Essay von \textsc{Schaefer} (Werke
XI, 2) befindet. — Ähnliche Einzelheiten liessen sich auch für andere Gebiete, so für die Astronomie und
die reine Mathematik, anführen.

Erst mit den beiden Essay-Bänden ist also die Herausgabe von \textsc{Gauss}’ Gesammelten Werken zu
dem befriedigenden Abschlusse gelangt, den \textsc{Klein} in dem ersten seiner (von 1898 bis 1920 in den Math.
Ann. veröffentlichten) Berichte über den Stand der Herausgabe von \textsc{Gauss}’ Werken auch insofern als erstrebenswertes
Ziel bezeichnet, als es sich um eine seit langer Zeit von der Göttinger Gesellschaft der
Wissenschaften übernommene Aufgabe handelt. Im besonderen geben die Essays eine Bestätigung dafür,
dass \textsc{Gauss}’ Tätigkeit, obwohl sie über ein halbes Jahrhundert zurückliegt, nicht einer vergangenen Periode
der Wissenschaft angehört, sondern in unmittelbarer Beziehung zu den Aufgaben der Gegenwart steht.
\begin{flushright}\textsc{M. Brendel.\quad L. Schlesinger.}\end{flushright}
\par\bigskip\centerline{\rule{.25\textwidth}{.3pt}}
\clearpage\SourceLeaf{698}
\begin{center}
{\Large\textls[120]{INHALT.}}\\[10pt]
{\large GAUSS WERKE BAND X\textsubscript{2}.}\\[10pt]
ABHANDLUNGEN\\
ÜBER GAUSS’ WISSENSCHAFTLICHE TÄTIGKEIT AUF DEN\\
GEBIETEN DER REINEN MATHEMATIK UND MECHANIK.
\end{center}
\newcommand{\FinalCRow}[4]{\par\noindent\hangindent=2em\hangafter=1 #1\enspace #2\hfill\mbox{#3}\nobreak\quad S.\ #4\par}
\small
\FinalCRow{1.}{\textsc{Paul Bachmann}: Über \textsc{Gauss}’ zahlentheoretische Arbeiten}{.\ .\ .\ .\ .\ .\ .\ .\ .\ .\ .\ .}{1—75}
\FinalCRow{2.}{\textsc{Ludwig Schlesinger}: Über \textsc{Gauss}’ Arbeiten zur Funktionentheorie}{.\ .\ .\ .\ .\ .\ .\ .\ .}{1—222}
\FinalCRow{3.}{\textsc{Alexander Ostrowski}: Über den ersten und vierten \textsc{Gauss}schen Beweis des Fundamentalsatzes der Algebra}{.\ .\ .\ .\ .\ .\ .\ .\ .\ .\ .\ .\ .\ .\ .\ .\ .\ .\ .\ .\ .\ .}{1—18}
\FinalCRow{4.}{\textsc{Paul Stäckel}: \textsc{Gauss} als Geometer}{.\ .\ .\ .\ .\ .\ .\ .\ .\ .\ .\ .\ .\ .\ .\ .\ .\ .\ .\ .}{1—123}
\FinalCRow{5.}{\textsc{Oskar Bolza}: Gauss und die Variationsrechnung}{.\ .\ .\ .\ .\ .\ .\ .\ .\ .\ .\ .\ .\ .\ .}{1—95}
\FinalCRow{6.}{\textsc{Philipp Maennchen}: Gauss als Zahlenrechner}{.\ .\ .\ .\ .\ .\ .\ .\ .\ .\ .\ .\ .\ .\ .\ .}{1—76}
\FinalCRow{7.}{\textsc{Harald Geppert}: Über \textsc{Gauss}’ Arbeiten zur Mechanik und Potentialtheorie}{.\ .\ .\ .\ .\ .}{1—61}
\par\bigskip\centerline{\rule{.25\textwidth}{.3pt}}
\clearpage\endgroup
\begingroup\NoSourceHead\SourceLeaf{701}
\setlength{\fboxsep}{7pt}
\noindent\fbox{\begin{minipage}{\dimexpr\textwidth-2\fboxsep-2\fboxrule\relax}
\fontsize{8}{9.4}\selectfont
\noindent\rule{\linewidth}{.2pt}
\begin{center}
{\large CARL FRIEDRICH GAUSS}\\[2pt]
{\Large WERKE}\\[3pt]
HERAUSGEGEBEN VON DER GESELLSCHAFT DER WISSENSCHAFTEN\\
ZU GÖTTINGEN\\[3pt]
IN KOMMISSION BEI JULIUS SPRINGER IN BERLIN\\[2pt]
\rule{3em}{.2pt}
\end{center}
\newcommand{\FPC}[2]{\par\noindent\hangindent=3em\hangafter=1\makebox[3em][r]{#1\enspace}#2\par}
\noindent Band
\FPC{I.}{DISQUISITIONES ARITHMETICAE. Zweiter Abdruck. 478 Seiten. 1870.\hfill Kart. RM 48.—}
\FPC{II.}{HÖHERE ARITHMETIK. Zweiter Abdruck. 528 Seiten. 1876.\hfill Kart. RM 53.—}
\FPC{}{NACHTRAG zum ersten Abdruck des zweiten Bandes. Seite 493—528. 1876.\hfill Kart. RM 3.80}
\FPC{III.}{ANALYSIS. Zweiter Abdruck. 499 Seiten. 1876.\hfill Kart. RM 50.—}
\FPC{IV.}{WAHRSCHEINLICHKEITSRECHNUNG UND GEOMETRIE. Zweiter Abdruck. 492 Seiten. 1880.\hfill Kart. RM 50.—}
\FPC{V.}{MATHEMATISCHE PHYSIK. Zweiter Abdruck. 642 Seiten. 1877.\hfill Kart. RM 64.—}
\FPC{VI.}{ASTRONOMISCHE ABHANDLUNGEN UND AUFSÄTZE. 664 Seiten. 1874.\hfill Kart. RM 69.—}
\FPC{VII.}{THEORIA MOTUS UND THEORETISCH-ASTRONOMISCHER NACHLASS. (Parabolische Bewegung, Störungen der Ceres und der Pallas, Theorie des Mondes.) 650 S. 1906.\hfill Kart. RM 65.—}
\FPC{VIII.}{ARITHMETIK. ANALYSIS, WAHRSCHEINLICHKEITSRECHNUNG, GEOMETRIE. (Nachträge zu Band I—IV.) 458 Seiten. 1900.\hfill Kart. RM 46.—}
\FPC{IX.}{GEODÄSIE. (Fortsetzung von Band IV.) 528 Seiten. 1903.\hfill Kart. RM 53.—}
\FPC{X.}{Abteilung 1. NACHLASS UND BRIEFWECHSEL ZUR REINEN MATHEMATIK. (Nachträge zu Band I—IV und VIII.) TAGEBUCH. 586 Seiten. 1917.\hfill Kart. RM 59.—}
\FPC{}{NACHBILDUNG DES TAGEBUCHS (Notizenjournals) 1796—1814.\hfill RM 1.20}
\FPC{X.}{Abteilung 2. ABHANDLUNGEN ÜBER GAUSS’ WISSENSCHAFTLICHE TÄTIGKEIT AUF DEN GEBIETEN DER REINEN MATHEMATIK UND MECHANIK. 672 Seiten. 1933.\hfill Brosch. RM 72.50.\quad Kart. RM 73.50}
\par\noindent\hspace*{5em}Hieraus einzeln:\par
\FPC{}{\hspace*{1em}Abhandlung 1 und 5: ÜBER GAUSS’ ZAHLENTHEORETISCHE ARBEITEN. Von \textsc{Paul Bachmann}. — GAUSS UND DIE VARIATIONSRECHNUNG. Von \textsc{Oscar Bolza}. \(75+95\) Seiten. 1922.\hfill Brosch. RM 17.—}
\FPC{}{\hspace*{1em}Abhandlung 2 und 3\quad ÜBER GAUSS’ ARBEITEN ZUR FUNKTIONENTHEORIE. Von \textsc{Ludwig Schlesinger}. — ÜBER DEN ERSTEN UND VIERTEN GAUSSSCHEN BEWEIS DES FUNDAMENTALSATZES DER ALGEBRA. Von \textsc{Alexander Ostrowski}. \(222+18\) Seiten. 1933.\hfill Brosch. RM 27.—}
\FPC{}{\hspace*{1em}Abhandlung 4\quad GAUSS ALS GEOMETER. Von \textsc{Paul Staeckel}. 123 Seiten. 1923.\hfill Brosch. RM 12.50}
\FPC{}{\hspace*{1em}Abhandlung 6. GAUSS ALS ZAHLENRECHNER. Von \textsc{Philipp Maennchen}. 76 Seiten. 1930.\hfill Brosch. RM 8.40}
\FPC{}{\hspace*{1em}Abhandlung 7\quad ÜBER GAUSS’ ARBEITEN ZUR MECHANIK UND POTENTIALTHEORIE. Von \textsc{Harald Geppert}. 63 Seiten. 1933.\hfill Brosch. RM 7.60}
\FPC{XI.}{Abteilung 1. NACHTRÄGE ZUR PHYSIK, CHRONOLOGIE UND ASTRONOMIE. (Nachträge zu Band V—VII.) 518 Seiten. 1927.\hfill Brosch. RM 47. —.\quad Kart. RM 48.—}
\FPC{XI.}{Abteilung 2. ABHANDLUNGEN ÜBER GAUSS’ WISSENSCHAFTLICHE TÄTIGKEIT AUF DEN GEBIETEN DER ANGEWANDTEN MATHEMATIK (GEODÄSIE, PHYSIK UND ASTRONOMIE). 640 Seiten. 1929.\hfill Brosch RM 70.—.\quad Kart. RM 71.—}
\par\noindent\hspace*{5em}Hieraus einzeln:\par
\FPC{}{\hspace*{1em}Abhandlung 1: ÜBER DIE GEODÄTISCHEN ARBEITEN VON GAUSS. Von \textsc{A. Galle}. 165 Seiten. 1924.\hfill Brosch. RM 17.—}
\FPC{}{\hspace*{1em}Abhandlung 2. ÜBER GAUSS’ PHYSIKALISCHE ARBEITEN (MAGNETISMUS, ELEKTRODYNAMIK, OPTIK). Von \textsc{Clemens Schaefer}. Mit 12 Figuren. 217 Seiten. 1929.\hfill Brosch. RM 24.—}
\FPC{}{\hspace*{1em}Abhandlung 3: ÜBER DIE ASTRONOMISCHEN ARBEITEN VON GAUSS. Von \textsc{Martin Brendel}. VI, 258 Seiten. 1929.\hfill Brosch. RM 29.—}
\FPC{XII.}{VARIA. ATLAS DES ERDMAGNETISMUS. Mit Textfiguren, 1 Übersichtstafel, 4 Zahlentafeln und 18 Kartentafeln. V, 415 Seiten. 1929.\hfill Brosch. RM 49.—.\quad Kart. RM 50.—}
\begin{center}\rule{3em}{.2pt}\end{center}
\par\noindent\hspace*{1em}\textit{Auf die Preise aller vor 1932 erschienenen Bände wird ein Notnachlass von 10\% gewährt.}
\par\medskip\noindent\rule{\linewidth}{.3pt}\par\vspace{-6pt}\noindent\rule{\linewidth}{.2pt}
\begin{center}\fontsize{6.5}{8}\selectfont GÖTTINGEN, GEDRUCKT IN DER DIETERICHSCHEN UNIVERSITÄTS-BUCHDRUCKEREI (W. FR. KAESTNER)\end{center}
\noindent\rule{\linewidth}{.15pt}
\end{minipage}}
\par\smallskip\centerline{\footnotesize PRINTED IN GERMANY.}
\clearpage\endgroup
